(1) 以平面饱和状态 \(p_{\rm sat}\) 为化学势基准，理想蒸气的摩尔每粒子化学势差为
\[
\mu_v(T,p)-\mu_v(T,p_{\rm sat})=k_BT\ln(p/p_{\rm sat}).
\]
曲率使液内 Laplace 压强高于气相压强 \(2\gamma/r\)。因液态近似不可压，
\(d\mu_\ell=v_\ell\,dp_\ell\)，所以小液滴相对平面共存液体的化学势增量（本模型忽略 \(v_\ell(p-p_{\rm sat})\) 的次级项）为 \(2\gamma v_\ell/r\)。平衡要求两者相等，故
\[
\boxed{p_{\rm eq}(r)=p_{\rm sat}\exp\!\left(\frac{2\gamma v_\ell}{r k_BT}\right)}.
\]
当 \(r\to\infty\) 恢复平面蒸气压；小半径要求更高蒸气分压才能平衡。因此在固定 \(S>1\) 下，太小的液滴会蒸发而非自发长大。这里的近似条件也要求蒸气压力对液体化学势的直接修正远小于 Laplace 主项。

(2) 每个由蒸气转入液态的分子使化学势降低
\(\Delta\mu=k_BT\ln S>0\)，体积驱动力 \(A=\Delta\mu/v_\ell\)。半径 \(r\) 的界面形成能为 \(4\pi\gamma r^2\)，体积收益为 \(n\Delta\mu=4\pi Ar^3/3\)。所以
\[
\boxed{\Delta G(r)=4\pi\gamma r^2-\frac{4\pi}{3}Ar^3},
\qquad A=\frac{k_BT\ln S}{v_\ell}.
\]
其一阶导数为 \(4\pi r(2\gamma-Ar)\)，非零驻点
\[
\boxed{r_*=\frac{2\gamma}{A}
=\frac{2\gamma v_\ell}{k_BT\ln S}},\quad
\boxed{n_*=\frac{4\pi r_*^3}{3v_\ell}}.
\]
二阶导数 \(8\pi\gamma-8\pi Ar\)，在临界点为 \(-8\pi\gamma<0\)。因此临界核是自由能势垒的极大值而非热力学稳定局部极小值。代回求得
\[
\boxed{\Delta G_*=\frac{4\pi}{3}\gamma r_*^2
=\frac{16\pi\gamma^3v_\ell^2}{3(k_BT\ln S)^2}}.
\]
对 \(0<r<r_*\) 的核，微小生长增加 \(\Delta G\)，易回缩；对 \(r>r_*\)，增长降低自由能，宏观液滴得以形成。

(3) Kelvin 平衡要求 \(\ln S=2\gamma v_\ell/(rk_BT)\)，所以其平衡半径正是 \(r_*\)；它是动力学意义的临界平衡，不代表稳定自由能谷。当 \(S\downarrow1\)，\(\ln S\to0^+\)，从而 \(r_*\sim(\ln S)^{-1}\)、\(n_*\sim(\ln S)^{-3}\)、\(\Delta G_*\sim(\ln S)^{-2}\to\infty\)。因此固定温度下成核率指数被强烈抑制。增大过饱和度使势垒快速下降，但当 \(r_*\) 接近分子尺度时连续表面张力、不可压液体与球形界面模型不再可靠，不能将 \(S\to\infty\) 的形式极限解释成精确微观结果。利用临界条件 \(n_*\Delta\mu=4\pi r_*^3A/3=(8\pi/3)\gamma r_*^2\)，得到
\[
\boxed{\Delta G_*=\frac12 n_*\Delta\mu}.
\]
进一步 \(\ln J=C-\Delta G_*/(k_BT)\)，若暂视前因子 \(C\) 不随 \(S\) 变化，则
\[
\frac{\partial\ln J}{\partial\ln S}
=\frac{2\Delta G_*}{k_BT\ln S}>0,
\]
展示弱过饱和区的指数级灵敏度。该表达式仅描述指数主项，不宣称真实动力学前因子保持常数。

(4) 将题给 Tolman 形式仅用于定义的界面自由能，得到
\[
\Delta G_{\mathrm T}(r)=4\pi\gamma(r^2-2\delta r)-\frac{4\pi}{3}Ar^3,\quad
\Delta G'_{\mathrm T}(r)=4\pi[2\gamma(r-\delta)-Ar^2].
\]
令 \(r_0=2\gamma/A\)，临界方程可写为
\(r^2-r_0r+r_0\delta=0\)。选择在 \(\delta\to0\) 时连续趋向 \(r_0\) 的大根，得
\[
r_{*,\mathrm T}=\frac{r_0}{2}\left(1+\sqrt{1-\frac{4\delta}{r_0}}\right)
=\boxed{r_0-\delta+O(\delta^2/r_0)}.
\]
对于 \(\delta>0\) 可能出现的另一根接近 \(\delta\)，不在 \(r\gg|\delta|\) 的模型自洽域，不能误认成可靠新物理稳定核。利用无修正势垒在 \(r_0\) 的一阶导数为零，代入临界半径的改变量在一阶不影响势垒，故
\[
\boxed{\Delta G_{*,\mathrm T}
=\Delta G_*-8\pi\gamma r_0\delta+O(\gamma\delta^2)}
\]
或 \(\Delta G_{*,\mathrm T}/\Delta G_*=1-6\delta/r_0+O((\delta/r_0)^2)\)。正 \(\delta\) 在本指定模型中降低临界半径及成核势垒，指数上提高成核倾向；但需要 \(|\delta|/r_0\ll1\)，更需液滴尺度远大于分子相关长度。切勿把线性近似外推到 \(r\sim\delta\) 或强过饱和的分子簇尺度。